\documentclass[11pt]{article}
\usepackage[margin=0.9in]{geometry}
\usepackage{amsmath,amssymb,amsthm}
\usepackage{graphicx}
\usepackage{float}
\usepackage{microtype}
\usepackage[hidelinks]{hyperref}

\newtheorem{theorem}{Theorem}
\newtheorem{lemma}{Lemma}
\newcommand{\Prob}{\mathcal P}
\newcommand{\TV}{\operatorname{TV}}
\newcommand{\C}{\mathbb C}

\title{A classical finite counterexample to the\\
$({\rm Lip}_p)$-isometric state-span conjecture}
\author{Candidate solution packet for run \texttt{fa\_banach\_001}}
\date{17 August 2026}

\begin{document}
\maketitle

\begin{center}
\fbox{\parbox{0.9\textwidth}{
\textbf{Status: candidate counterexample, likely valid.}
For every $p\geq1$, the regular action of the classical group $C_3$ on a
three-point path metric has
$\overline{\operatorname{span}}^{w^*}S_p(A)=A^*$, although $S_p(A)$ is a
proper subset of the state space.  This disproves Conjecture~4.13 of
Chirvasitu.
}}
\end{center}

\section{Source conjecture}

For a compact quantum group $A$ acting on a compact metric space $(X,d)$,
Chirvasitu defines a state $\psi\in S(A)$ to be
$({\rm Lip}_p)$-isometric when
\[
 W_p(x\mathbin{\triangleleft}\psi,y\mathbin{\triangleleft}\psi)
 \leq d(x,y)\qquad(x,y\in X).
\]
The set of such states is denoted by $S_p(A)$.  Conjecture~4.13 on PDF
page~23 of arXiv:1407.0960 \cite{source} states:

\begin{quote}
Let $A$ act faithfully on the compact metric space $(X,d)$, and $p\geq1$.
The set of states in the weak$^*$-closed linear span of $S_p(A)$ coincides
with $S_p(A)$.
\end{quote}

\begin{figure}[H]
\centering
\includegraphics[width=0.97\textwidth]{figures/open_problem_crop.png}
\caption{Definition~4.12 and Conjecture~4.13 in the source (PDF page~23).}
\end{figure}

\section{Idea of the proof}

Classical states are probability measures on the acting group.  For a
transitive finite action, Haar measure sends every point of $X$ to the same
uniform distribution.  Measures sufficiently close to Haar therefore send
every pair of points to distributions whose Wasserstein distance is smaller
than the least nonzero distance in $X$.  Thus $S_p(A)$ contains a relative
neighborhood of Haar and consequently spans the whole dual $A^*$.

This is incompatible with any non-isometric element of the group: its point
mass is a state of $A$ and hence lies in that full span, but it expands some
distance and so is not in $S_p(A)$.  The three-cycle acting on a three-point
path is the smallest transparent example.

\section{A finite transport estimate}

We use $\TV(\alpha,\beta)=\frac12\|\alpha-\beta\|_1$ for probability
measures on a finite set.

\begin{lemma}\label{lem:transport}
Let $(X,d)$ be a finite metric space of diameter $D$.  For all
$\alpha,\beta\in\Prob(X)$ and $p\geq1$,
\[
 W_p(\alpha,\beta)^p\leq D^p\TV(\alpha,\beta).
\]
\end{lemma}

\begin{proof}
Put the common mass $\min\{\alpha(z),\beta(z)\}$ on $(z,z)$ for every
$z\in X$.  The unmatched mass on either side has total mass
$r=\TV(\alpha,\beta)$.  Couple those two residual measures arbitrarily.
The diagonal part costs zero and every residual pair costs at most $D^p$,
so the resulting transport plan has cost at most $D^p r$.
\end{proof}

\section{A general counterexample principle}

\begin{theorem}\label{thm:general}
Let a finite group $G$ act faithfully and transitively on a finite metric
space $(X,d)$.  Suppose the action is not isometric.  Put $A=C(G)$.  Then,
for every $p\geq1$,
\[
 \overline{\operatorname{span}}^{w^*}S_p(A)=A^*,
 \qquad S_p(A)\ne S(A).
\]
In particular, the states in the weak$^*$-closed linear span of $S_p(A)$ do
not coincide with $S_p(A)$.
\end{theorem}

\begin{proof}
A state of $A=C(G)$ is a probability measure $\mu$ on $G$, and
$x\mathbin{\triangleleft}\mu$ is the pushforward of $\mu$ under
$g\mapsto xg$.  Let $h$ be the uniform probability measure on $G$.  By
transitivity, $x\mathbin{\triangleleft}h$ is the uniform probability measure
$u_X$ on $X$, independently of $x$.

Let $D=\operatorname{diam}X$ and let
$m=\min\{d(x,y):x\ne y\}>0$.  Pushforward contracts total variation, hence
for every $x,y\in X$,
\begin{align*}
 \TV(x\mathbin{\triangleleft}\mu,
       y\mathbin{\triangleleft}\mu)
 &\leq \TV(x\mathbin{\triangleleft}\mu,u_X)
       +\TV(u_X,y\mathbin{\triangleleft}\mu)\\
 &\leq 2\TV(\mu,h).
\end{align*}
Lemma~\ref{lem:transport} gives
\[
 W_p(x\mathbin{\triangleleft}\mu,
       y\mathbin{\triangleleft}\mu)^p
 \leq 2D^p\TV(\mu,h).
\]
Thus every state satisfying
$\TV(\mu,h)<m^p/(2D^p)$ belongs to $S_p(A)$: the case $x=y$ is trivial,
and otherwise the displayed Wasserstein distance is less than
$m\leq d(x,y)$.

Consequently $S_p(A)$ contains a relative neighborhood of the interior point
$h$ in the probability simplex on $G$.  Its real affine span is the whole
mass-one hyperplane, and its real linear span is the space of all signed
measures on $G$.  After complexification, its complex linear span is $A^*$.
Since $A^*$ is finite-dimensional, weak$^*$ closure changes nothing.

It remains to see that $S_p(A)$ is proper.  Since the action is not
isometric, some group element expands a distance.  Indeed, if an element
strictly contracts some pair, its inverse expands the image pair; hence there
are $g\in G$ and $x,y\in X$ with
$d(xg,yg)>d(x,y)$.  For the state $\delta_g$,
\[
 W_p(x\mathbin{\triangleleft}\delta_g,
       y\mathbin{\triangleleft}\delta_g)
 =W_p(\delta_{xg},\delta_{yg})=d(xg,yg)>d(x,y).
\]
Therefore $\delta_g\notin S_p(A)$, although it lies in the linear span just
proved to be all of $A^*$.
\end{proof}

\section{The explicit three-point example}

Take $X=\{0,1,2\}$ with
\[
 d(0,1)=d(1,2)=1,\qquad d(0,2)=2,
\]
and let $G=C_3=\mathbb Z/3\mathbb Z$ act by
$xg=x+g\pmod 3$.  This regular action is faithful and transitive.  It is not
isometric because addition by one sends the pair $(1,2)$, at distance one,
to $(2,0)$, at distance two.  Faithfulness in the source's coaction sense is
also immediate: the regular action has trivial kernel (equivalently, its
matrix coefficients contain all point masses in $C(C_3)$).  Theorem~\ref{thm:general}
therefore disproves Conjecture~4.13 for every $p\geq1$.

For completeness, the full-span claim can be witnessed by only three states.
Let $h=(1/3,1/3,1/3)$ and choose
$0<\varepsilon<2^{-(p+1)}$.  Then
\[
 \mu_0=h,\qquad
 \mu_1=h+\varepsilon(\delta_0-\delta_2),\qquad
 \mu_2=h+\varepsilon(\delta_1-\delta_2)
\]
are states in $S_p(A)$, because $D=2$, $m=1$, and
$\TV(\mu_i,h)\leq\varepsilon$.  As row vectors, their determinant is
$\varepsilon^2\ne0$, so they span $A^*\cong\C^3$.  Yet
\[
 W_p(1\mathbin{\triangleleft}\delta_1,
       2\mathbin{\triangleleft}\delta_1)
 =W_p(\delta_2,\delta_0)=2>1=d(1,2),
\]
so $\delta_1\notin S_p(A)$.

\section{Verification and novelty status}

The included script \texttt{code/verify\_c3\_transport.py} uses exact rational
arithmetic and a finite min-cost-flow solver to compute all nine orbit-law
transport problems for the three displayed spanning states at
$p=1,2,3,4,5,8$.  It verifies their determinant and the point-mass violation.
The calculation is a regression check only; Lemma~\ref{lem:transport} and the
proof above establish the result for every real $p\geq1$.

A bounded novelty check on 17 August 2026 searched arXiv, exact web phrases,
the visible citation record, and later same-author quantum-isometry papers for
the arXiv id, paper title, ``Conjecture 4.13,'' ``$S_p(A)$,''
``weak$^*$-closed linear span,'' and close variants involving
$({\rm Lip}_p)$-isometric states.  No later proof or counterexample to this
span conjecture was found.  Novelty is plausible, not certified.

\section{Scope and human-review recommendation}

This packet fully negates Conjecture~4.13, even for a classical finite group.
It does not address the separate conjecture earlier in the source that the
conjunction of all $({\rm Lip}_p)$ conditions implies $(D)$-isometry.
Human review should first check the total-variation neighborhood argument and
the classical identification of states with group measures.  All remaining
claims reduce to the displayed three-point metric and elementary linear
algebra.

\begin{thebibliography}{9}
\bibitem{source}
A.~Chirvasitu,
\emph{On quantum symmetries of compact metric spaces},
J. Geom. Phys. 87 (2015), 134--151;
\href{https://arxiv.org/abs/1407.0960}{arXiv:1407.0960}.
\end{thebibliography}

\end{document}
